\section{完成定义：拒稿意见逐条映射}
\label{sec:dod}

本计划的验收不以``跑完了多少实验''衡量，而以下表逐条对照：拒稿意见的每个具体构成，在哪个实验完成后、以论文中的哪类证据形式，被认为已回应。此表同时是修订说明信（response letter，若重投或改投时需要）的骨架。

\begin{table}[htbp]
\centering
\caption{拒稿意见构成与补强证据的映射。``章节''指 paper-journal-v2 的对应位置。}
\label{tab:dod}
\small
\begin{tabularx}{\textwidth}{@{}p{4.2cm}XX@{}}
\toprule
拒稿构成 & 补强实验 & 完成形态 \\
\midrule
新颖性不足（第一阶段） & 文献审计 + 叙事重构（已完成的 v2） & 新标题与 formulation 摘要；相关工作七维设计空间表与 delta 段；贡献列表重排为``形式化$\rightarrow$方法论$\rightarrow$发现$\rightarrow$实例'' \\
端到端规模（$n{=}25$、单点） & A0+A1 & $n{=}300$ 端到端主表，区间半宽 $\approx\pm3.3$ F1 点；BM25 进入端到端层 \\
跨数据集端到端缺失 & A2 & HotpotQA 与 2Wiki 各 $n{=}100$ 端到端表（视资源） \\
单阅读器依赖 & B1 & 三阅读器趋势稳定性表/图；``对某阅读器成立''限定词删除或改写 \\
位置剖面借用 & B2 & 自测阅读器侧位置曲线图，替换借用的动机段 \\
消融未成矩阵 & C1+C2+C3 & 三数据集消融表；配置$\times$预算矩阵表；升级版三轴 frontier 图 \\
召回--F1 相关 $n$ 小 & A0 扩容 + B3 & 相关研究更新为 $n{=}100$ 起步、跨阅读器分组 \\
``More work needed''整体 & 全部 + 本两阶段文档 & 修订说明信按本表逐条给出``意见$\rightarrow$改动$\rightarrow$证据''三段式回应 \\
\bottomrule
\end{tabularx}
\end{table}

最后给出三档验收口径，供按实际资源停机。最低档（P0 完成：A0 + 第一阶段 v2）：端到端 $n{=}100$ 落盘、叙事重构完成——此时``premature''的第一构成已实质缩小但未完全消解，适合先行内部评审。标准档（P0+P1：A1 + C1--C3）：矩阵化证据齐备、端到端 $n{=}300$——达到可重投或改投（KBS/ESWA 类）的证据门槛，且两条 premature 构成全部消解。完整档（加 P2）：跨数据集端到端与三阅读器稳健性——使论文对``结构性权衡''的声明具备跨阅读器与跨任务形态的证据面，投稿竞争力最高。三档之间随时可停，已完成的每一格都独立可写——这是本计划与``盲目扩容''的根本区别。
